\documentclass[11pt]{article}
\usepackage{mathblatt}
\begin{document}
\weit
\typeout{PROBE-BEGIN ableitung-und-aenderungsrate-e3-k1-s1-v1}
\begin{aufgabe}{\detokenize{ableitung-und-aenderungsrate-e3-k1-s1-v1}}
\begin{teile}
\teil Probe

\wertetabelle{h}{m}{1,0{,}5,0{,}1,0{,}01}
\end{teile}
\end{aufgabe}
\typeout{PROBE-END ableitung-und-aenderungsrate-e3-k1-s1-v1}
\typeout{PROBE-BEGIN extremalprobleme-e1-k2-s3-v3}
\begin{aufgabe}{\detokenize{extremalprobleme-e1-k2-s3-v3}}
\begin{teile}
\teil Probe

\wertetabelle{u}{A(u)}{1,2,3,4,5}
\end{teile}
\end{aufgabe}
\typeout{PROBE-END extremalprobleme-e1-k2-s3-v3}
\typeout{PROBE-BEGIN funktionsklassen-und-eigenschaften-e6-k1-s1-v1}
\begin{aufgabe}{\detokenize{funktionsklassen-und-eigenschaften-e6-k1-s1-v1}}
\begin{teile}
\teil Probe

\wertetabelle{x}{f(x)}{-2,-1,0,1,2}\begin{ksys}[xmin=-3,xmax=3,ymin=-3,ymax=3]\end{ksys}
\end{teile}
\end{aufgabe}
\typeout{PROBE-END funktionsklassen-und-eigenschaften-e6-k1-s1-v1}
\typeout{PROBE-BEGIN kreis-e1-k6-s4-v1}
\begin{aufgabe}{\detokenize{kreis-e1-k6-s4-v1}}
\begin{teile}
\teil Probe

\wertetabelle{$d$ in cm}{$u$ in cm}{1,2,3,6}
\end{teile}
\end{aufgabe}
\typeout{PROBE-END kreis-e1-k6-s4-v1}
\typeout{PROBE-BEGIN lineare-funktionen-e1-k1-s0-v1}
\begin{aufgabe}{\detokenize{lineare-funktionen-e1-k1-s0-v1}}
\begin{teile}
\teil Probe

\wertetabelle[3,6,9,12]{x}{y}{1,2,3,4}
\end{teile}
\end{aufgabe}
\typeout{PROBE-END lineare-funktionen-e1-k1-s0-v1}
\end{document}
